% Original: PDF strana 21, gornji slajd.
\slajdid{02-s041}
\begin{frame}[label=02-s041]{Algoritam minimizacije}
\setlength{\parskip}{0pt}
% element: 02-s041:e01
Iz funkcionalne tabele ili specifikacije kombinacione mreže unosimo vrednosti funkcije u odgovarajuća polja Karnoove karte, prema stanjima ulaznih promenljivih.
\par\vspace{1.5mm}
% element: 02-s041:e02
\begin{enumerate}
\item Za funkciju u obliku \textbf{zbira proizvoda}, sve logičke jedinice u karti treba prekriti \alert{što manjim brojem površina što višeg reda}.
\item Za funkciju u obliku \textbf{proizvoda zbirova}, sve logičke nule u karti treba prekriti \alert{što manjim brojem površina što višeg reda}.
\end{enumerate}
\par\vspace{2mm}
% element: 02-s041:e03
\Oznaka{Uslov minimalnosti}\par\vspace{1mm}
Ako je broj površina minimalan, a površine najveće moguće, algoritam garantuje minimalnu funkciju. Pod tim uslovom nije moguća jednostavnija realizacija pomoću I i ILI logičkih kola.
\par\vspace{1.5mm}
\Rezultat{Podrazumevamo da su na ulazu kombinacione mreže dostupne i prave i komplementne vrednosti signala.}
\end{frame}
